\addcontentsline{toc}{subsubsection}{Solution 19. The Reciprocal-Slope Pinch}

\begin{dlsummary}[title={\bfseries Solution 19. The Reciprocal-Slope Pinch},breakable]

\begin{proof}
\textbf{Step 1. Establish a lower bound.}

Since $f(xf(x))>0$, the given inequality implies
\[
f(x+y)>y\qquad(x,y>0).
\]
Fix $z>0$ and take any $x\in(0,z)$ with $y=z-x$. Then
$f(z)>z-x$. As this holds for every such $x$, it follows that
\[
f(z)\geq z\qquad(z>0). \tag{1}
\]

\medskip
\noindent
\textbf{Step 2. Control the excess over the identity.}

Define $g(x)=f(x)-x$, so $g(x)\geq0$ by (1). Since
$f(xf(x))=xf(x)+g(xf(x))$, the original inequality is equivalent to
\[
g(x+y)-g(x)\geq\frac{g(xf(x))}{x}\qquad(x,y>0). \tag{2}
\]
In particular, $g$ is nondecreasing.

If $f(1)>1$, taking $x=1$ and $y=f(1)-1$ in the original inequality
would give
\[
f(1)-1+f(f(1))\leq f(f(1)),
\]
which is impossible. Hence $f(1)\leq1$, and (1) gives $f(1)=1$;
therefore $g(1)=0$. Since $g$ is nonnegative and nondecreasing,
\[
g(x)=0\qquad(0<x\leq1). \tag{3}
\]

\medskip
\noindent
\textbf{Step 3. Rule out positive excess above $1$.}

Suppose $g(a)=c>0$ for some $a>1$. For every $u>1$, we have
$uf(u)\geq u^2>u$. Monotonicity of $g$ and (2) therefore imply, for
all $v>u>1$,
\[
g(v)\geq g(u)+\frac{g(uf(u))}{u}
\geq g(u)+\frac{g(u)}{u}. \tag{4}
\]
Fix $X>a$. For any positive integer $N$, choose
\[
u_j=a+\frac{j(X-a)}{N+1}\qquad(0\leq j\leq N).
\]
All these points lie in $[a,X)$, so $g(u_j)\geq c$ and
$u_j<X$. Applying (4) to each successive pair gives
\[
g(u_{j+1})\geq g(u_j)+\frac{c}{X}
\qquad(0\leq j<N).
\]
Thus $g(u_N)\geq c+Nc/X$. But monotonicity gives $g(u_N)\leq g(X)$,
which is a fixed finite value. Choosing $N$ sufficiently large is a
contradiction. Hence $g(x)=0$ for all $x>1$ as well. Together with (3),
this proves $f(x)=x$ for every $x>0$.

Finally, the identity function satisfies the original inequality with
equality throughout. It is therefore the unique solution.
\end{proof}

\end{dlsummary}